\documentclass[12pt,a4paper]{ctexart}

\usepackage{amsmath,amssymb}
\usepackage{geometry}
\usepackage{enumitem}
\usepackage{hyperref}
\usepackage{booktabs}
\usepackage{listings}

\geometry{left=2.5cm,right=2.5cm,top=2.5cm,bottom=2.5cm}
\setlist{nosep,leftmargin=2em}

\title{第3章\quad 功和能}
\author{}
\date{}

\begin{document}

\maketitle

\section{知识点总结}

\subsection{功}

\begin{itemize}
    \item \textbf{恒力的功}：$A = Fs\cos\theta = \vec{F}\cdot\vec{S}$。
    \item \textbf{变力的功}：元功 $\mathrm{d}A = \vec{F}\cdot\mathrm{d}\vec{r}$；沿路径 $L$ 从 $a$ 到 $b$：$A = \displaystyle\int_{a(L)}^b \vec{F}\cdot\mathrm{d}\vec{r}$。
    \begin{itemize}
        \item 直角坐标：$A = \displaystyle\int_{a(L)}^b (F_x\,\mathrm{d}x + F_y\,\mathrm{d}y + F_z\,\mathrm{d}z)$。
        \item 自然坐标：$A = \displaystyle\int_{a(L)}^b F\cos\theta\,\mathrm{d}s$。
    \end{itemize}
    \item \textbf{说明}：
    \begin{enumerate}
        \item 功是标量，可正可负；
        \item 合力的功等于各分力功的代数和；
        \item 一般功与路径有关。
    \end{enumerate}
    解题第一步：写出 $\mathrm{d}A$，把 $\mathrm{d}A$ 写成"正在变化的那个变量"的函数。
    \item \textbf{功率}：
    \begin{itemize}
        \item 平均功率 $\bar P = \dfrac{\Delta A}{\Delta t}$；
        \item 瞬时功率 $P = \dfrac{\mathrm{d}A}{\mathrm{d}t} = \vec{F}\cdot\vec{v} = Fv\cos\theta$。
    \end{itemize}
\end{itemize}

\subsection{几种常见力的功}

\begin{itemize}
    \item \textbf{重力的功}（只与始末高度差有关）：$A_{\text{重}} = mg(z_1 - z_2)$，路径无关。
    \begin{itemize}
        \item 书上例 3.2：长 $L$ 质量 $M$ 的柔绳，B 端被提到与 A 端同高时重力作功 $A = -\frac{1}{4}MgL$（利用提起的部分质心 $y/2$）。
    \end{itemize}
    \item \textbf{弹性力的功}（胡克定律 $F = -kx$）：$A = \displaystyle\int_{x_1}^{x_2} -kx\,\mathrm{d}x = \dfrac{1}{2}kx_1^2 - \dfrac{1}{2}kx_2^2$。只与始末形变有关；变形减小作正功。
    \item \textbf{万有引力的功}：$A = \displaystyle\int_{r_1}^{r_2} -G\dfrac{mM}{r^2}\mathrm{d}r = GmM\left(\dfrac{1}{r_2} - \dfrac{1}{r_1}\right)$。移近时作正功，远离时作负功。
    \item \textbf{摩擦力的功}：方向始终与速度方向相反，$A = -\mu mgs$，与路径有关（非保守力）。
\end{itemize}

\subsection{动能定理}

\begin{itemize}
    \item \textbf{质点动能定理}：$A = \dfrac{1}{2}mv_2^2 - \dfrac{1}{2}mv_1^2 = E_{k2} - E_{k1}$。
    \begin{itemize}
        \item 功是过程量，动能是状态量。
        \item 只适用于惯性系。
    \end{itemize}
    \item \textbf{质点系动能定理}：$\displaystyle\sum A_i = \sum A_{i\text{外}} + \sum A_{i\text{内}} = \sum \dfrac{1}{2}m_iv_{i2}^2 - \sum \dfrac{1}{2}m_iv_{i1}^2$。
    \begin{itemize}
        \item 内力的矢量和为零，但内力功的代数和不一定为零（如炸弹爆炸）。
        \item 整链条类问题中各部分间内力功之和为零。
    \end{itemize}
\end{itemize}

\subsection{势能\quad 机械能守恒定律}

\begin{itemize}
    \item \textbf{保守力}：力所作的功与路径无关，只取决于始末相对位置；$\displaystyle\oint_L \vec{f}\cdot\mathrm{d}\vec{r} = 0$。重力、万有引力、弹性力都是保守力；摩擦力是非保守力。
    \item \textbf{势能}：相对量，无绝对大小；零点任意。
    \begin{itemize}
        \item 重力势能：$E_p = mgz$（零点取在 $z=0$ 平面）。
        \item 弹性势能：$E_p = \dfrac{1}{2}kx^2$（零点取在弹簧原长）。
        \item 万有引力势能：$E_p = -G\dfrac{mM}{r}$（零点取在 $r\to\infty$）。
        \item 球内引力势能 $E_p = -GMm\dfrac{3R^2 - x^2}{2R^3}$。
    \end{itemize}
    \item \textbf{保守力与势能关系}：
    \begin{itemize}
        \item 功是势能增量的负值：$A = -(E_{p2} - E_{p1}) = -\Delta E_p$。
        \item 保守力是势能的负梯度：$\vec{F} = -\nabla E_p$，$\nabla = \dfrac{\partial}{\partial x}\vec{i} + \dfrac{\partial}{\partial y}\vec{j} + \dfrac{\partial}{\partial z}\vec{k}$。
        \item 二维判断保守力：$\dfrac{\partial F_x}{\partial y} = \dfrac{\partial F_y}{\partial x}$。
    \end{itemize}
    \item \textbf{平衡稳定性}：
    \begin{itemize}
        \item 稳定平衡 $\dfrac{\mathrm{d}^2 E_p}{\mathrm{d}x^2} > 0$；
        \item 不稳定平衡 $\dfrac{\mathrm{d}^2 E_p}{\mathrm{d}x^2} < 0$；
        \item 随遇平衡 $\dfrac{\mathrm{d}^2 E_p}{\mathrm{d}x^2} = 0$。
    \end{itemize}
    \item \textbf{机械能守恒定律}：当 $A_{\text{外}} + A_{\text{非内}} = 0$ 时，$\Delta E = 0$，即 $E = E_k + E_p = \text{常数}$。
    \begin{itemize}
        \item 是对系统而言、对整个过程而言。
    \end{itemize}
\end{itemize}

\subsection{能量守恒定律}

\begin{itemize}
    \item 能量不能消失也不能创造，只能从一种形式转换为另一种形式；封闭系统的总能量守恒。
    \item 功是能量交换或转换的一种度量；机械能守恒是普遍能量守恒定律的特例。
\end{itemize}

\section{公式汇总}

\begin{enumerate}[leftmargin=3em,label=\textbf{公式\arabic*}：,ref=公式\arabic*]

    \item $A = \displaystyle\int_{a(L)}^b \vec{F}\cdot\mathrm{d}\vec{r} = \int (F_x\,\mathrm{d}x + F_y\,\mathrm{d}y + F_z\,\mathrm{d}z)$

    说明：变力沿路径的功（线积分），$\mathrm{d}\vec{r}$ 是位移元；直角坐标下三个分量的功的代数和。

    \item $A = \displaystyle\int_{a(L)}^b F\cos\theta\,\mathrm{d}s$

    说明：自然坐标下的变力功表示，$F\cos\theta$ 是力沿切向方向的分量，$\mathrm{d}s$ 是弧长元。

    \item $P = \vec{F}\cdot\vec{v} = Fv\cos\theta$

    说明：瞬时功率，$P$ 的单位为 W（瓦特），1 W = 1 J/s。

    \item $A_{\text{重}} = mg(z_1 - z_2)$

    说明：重力做的功等于重力大小乘以始末位置的高度差；只与始末位置有关，路径无关。

    \item $A_{\text{弹}} = \dfrac{1}{2}kx_1^2 - \dfrac{1}{2}kx_2^2$

    说明：弹性力做的功，$x_1,x_2$ 为弹簧形变量；变形减小时作正功。

    \item $A_{\text{引}} = GmM\left(\dfrac{1}{r_2} - \dfrac{1}{r_1}\right)$

    说明：万有引力做的功，$r_1,r_2$ 为质点到引力中心的距离；引力为吸引力，远离作负功。

    \item $A_{\text{摩}} = -\mu mgs$

    说明：滑动摩擦力做的功，$s$ 是路径长度（注意是路径不是位移），与路径有关。

    \item $A = E_{k2} - E_{k1} = \dfrac{1}{2}mv_2^2 - \dfrac{1}{2}mv_1^2$

    说明：质点动能定理，合力在某一路程上对质点做的功等于质点始末状态动能的增量；功是过程量，动能是状态量。

    \item $\displaystyle\oint_L \vec{f}\cdot\mathrm{d}\vec{r} = 0$

    说明：保守力沿任意闭合路径一周做的功为零；这是保守力等价判据。

    \item $E_{p\text{重}} = mgz$

    说明：重力势能（取 $z=0$ 为零点），与水平位置 $x,y$ 无关；势能零点的选取是任意的。

    \item $E_{p\text{弹}} = \dfrac{1}{2}kx^2$

    说明：弹性势能，零点取在弹簧原长（$x=0$）；$k$ 为劲度系数。

    \item $E_{p\text{引}} = -G\dfrac{mM}{r}$

    说明：万有引力势能，零点取在 $r\to\infty$；负号表示束缚态。

    \item $A_{\text{保}} = -(E_{p2} - E_{p1}) = -\Delta E_p$

    说明：保守力做的功等于势能增量的负值。

    \item $\vec{F} = -\nabla E_p = -\left(\dfrac{\partial E_p}{\partial x}\vec{i} + \dfrac{\partial E_p}{\partial y}\vec{j} + \dfrac{\partial E_p}{\partial z}\vec{k}\right)$

    说明：保守力是势能梯度的负值；势能曲线上某点斜率的负值就是该点处的保守力。

    \item $E_k + E_p = \text{常数}$（条件 $A_{\text{外}} + A_{\text{非内}} = 0$）

    说明：机械能守恒定律，机械能为动能与势能之和。

\end{enumerate}

\end{document}
